% Part: sets-functions-relations
% Chapter: sets
% Section: important-sets

\documentclass[../../../include/open-logic-section]{subfiles}

\begin{document}

\olfileid{sfr}{set}{imp}
\olsection{काही महत्त्वाचे संच}

\begin{ex}
आपण प्रामुख्याने ज्यातील !!{element}s या गणिती वस्तू आहेत अशा संचांचा
विचार करणार आहोत. असे चार संच इतके महत्त्वाचे आहेत की त्यांना
विशिष्ट नावे दिलेली आहेत:
\begin{multline*}
    \Nat = \{0, 1, 2, 3, \ldots\} \\
    \shoveright{\text{नैसर्गिक संख्यांचा संच}}\\
    \shoveleft{\Int = \{\ldots, -2, -1, 0, 1, 2, \ldots\}} \\
    \shoveright{\text{पूर्णांकांचा संच}}\\
    \shoveleft{\Rat = \Setabs{\nicefrac{m}{n}}{m, n \in \Int\text{ आणि }n \neq 0}}\\
    \shoveright{\text{परिमेय संख्यांचा संच}}\\
    \shoveleft{\Real = (-\infty, \infty)}\\
    \text{वास्तव संख्यांचा संच (सांतत्य)}
\end{multline*}
हे सर्व \emph{अनंत} संच आहेत; म्हणजेच, प्रत्येकात अनंत !!{element}s आहेत.

या संचांतून पुढे जाताना आपण आपल्या संग्रहात \emph{आणखी} संख्या समाविष्ट
करत असतो. खरे तर, $\Nat \subseteq \Int
\subseteq \Rat \subseteq \Real$ हे स्पष्ट असावे: कारण प्रत्येक नैसर्गिक संख्या
पूर्णांक असते; प्रत्येक पूर्णांक परिमेय असतो; आणि प्रत्येक परिमेय संख्या वास्तव असते.
त्याचप्रमाणे, $\Nat \subsetneq \Int \subsetneq
\Rat$ हेही स्पष्ट असावे, कारण $-1$ हा पूर्णांक आहे पण नैसर्गिक संख्या नाही,
आणि $\nicefrac{1}{2}$ ही परिमेय संख्या आहे पण पूर्णांक नाही.
$\Rat \subsetneq \Real$, म्हणजे परिमेय नसलेल्या काही वास्तव संख्या आहेत,
ही गोष्ट तितकी सहज स्पष्ट होत नाही\oliflabeldef{sfr:arith:real:realline}{;
पण आपण \olref[arith][real]{realline} मध्ये पुन्हा तिचा विचार करू}{}.

कधीकधी आपण धन पूर्णांकांचा संच $\PosInt = \{1,
2, 3, \dots\}$ आणि फक्त पहिल्या दोन नैसर्गिक संख्या असणारा संच
$\Bin = \{0, 1\}$ यांचाही उपयोग करू.
\end{ex}

\begin{tagblock}{compsci}
\begin{ex}[चिन्हमाला]
आणखी एक लक्षवेधी उदाहरण म्हणजे $A$ या वर्णमालेवरील \emph{सांत चिन्हमालांचा}
संच $A^{*}$: $A$ मधील घटकांचा कोणताही सांत अनुक्रम हा $A$ वरील चिन्हमाला
असतो. प्रत्येक $A$ वर्णमालेसाठी, $A$ वरील चिन्हमालांमध्ये आपण
\emph{रिक्त चिन्हमाला $\Lambda$} हिचाही समावेश करतो. उदाहरणार्थ,
\begin{multline*}
\Bin^*
=\{\Lambda,0,1,00,01,10,11,\\
000,001,010,011,100,101,110,111,0000,\ldots\}.
\end{multline*}
$x=x_{1}\ldots x_{n}\in A^{*}$ ही $A$ मधील $n$ ``अक्षरांनी'' बनलेली
चिन्हमाला असेल, तर तिची \emph{लांबी} $n$ आहे असे म्हणतो आणि $\len{x}=n$ असे लिहितो.
\end{ex}
\end{tagblock}

\begin{ex}[अनंत अनुक्रम]
कोणत्याही $A$ संचासाठी आपण $A$ मधील !!{element}s यांच्या अनंत अनुक्रमांचा
संच $A^\omega$ याचाही विचार करू शकतो. अनंत अनुक्रम
$a_1a_2a_3a_4\dots$ म्हणजे वस्तूंची एका दिशेने अनंत लांब जाणारी यादी;
यादीतील प्रत्येक वस्तू $A$ चा !!a{element} असते.
\end{ex}

\end{document}
